% Options for packages loaded elsewhere
% Options for packages loaded elsewhere
\PassOptionsToPackage{unicode}{hyperref}
\PassOptionsToPackage{hyphens}{url}
%
\documentclass[
  letterpaper,
]{scrbook}
\usepackage{xcolor}
\usepackage{amsmath,amssymb}
\setcounter{secnumdepth}{5}
\usepackage{iftex}
\ifPDFTeX
  \usepackage[T1]{fontenc}
  \usepackage[utf8]{inputenc}
  \usepackage{textcomp} % provide euro and other symbols
\else % if luatex or xetex
  \usepackage{unicode-math} % this also loads fontspec
  \defaultfontfeatures{Scale=MatchLowercase}
  \defaultfontfeatures[\rmfamily]{Ligatures=TeX,Scale=1}
\fi
\usepackage{lmodern}
\ifPDFTeX\else
  % xetex/luatex font selection
\fi
% Use upquote if available, for straight quotes in verbatim environments
\IfFileExists{upquote.sty}{\usepackage{upquote}}{}
\IfFileExists{microtype.sty}{% use microtype if available
  \usepackage[]{microtype}
  \UseMicrotypeSet[protrusion]{basicmath} % disable protrusion for tt fonts
}{}
\makeatletter
\@ifundefined{KOMAClassName}{% if non-KOMA class
  \IfFileExists{parskip.sty}{%
    \usepackage{parskip}
  }{% else
    \setlength{\parindent}{0pt}
    \setlength{\parskip}{6pt plus 2pt minus 1pt}}
}{% if KOMA class
  \KOMAoptions{parskip=half}}
\makeatother
% Make \paragraph and \subparagraph free-standing
\makeatletter
\ifx\paragraph\undefined\else
  \let\oldparagraph\paragraph
  \renewcommand{\paragraph}{
    \@ifstar
      \xxxParagraphStar
      \xxxParagraphNoStar
  }
  \newcommand{\xxxParagraphStar}[1]{\oldparagraph*{#1}\mbox{}}
  \newcommand{\xxxParagraphNoStar}[1]{\oldparagraph{#1}\mbox{}}
\fi
\ifx\subparagraph\undefined\else
  \let\oldsubparagraph\subparagraph
  \renewcommand{\subparagraph}{
    \@ifstar
      \xxxSubParagraphStar
      \xxxSubParagraphNoStar
  }
  \newcommand{\xxxSubParagraphStar}[1]{\oldsubparagraph*{#1}\mbox{}}
  \newcommand{\xxxSubParagraphNoStar}[1]{\oldsubparagraph{#1}\mbox{}}
\fi
\makeatother

\usepackage{color}
\usepackage{fancyvrb}
\newcommand{\VerbBar}{|}
\newcommand{\VERB}{\Verb[commandchars=\\\{\}]}
\DefineVerbatimEnvironment{Highlighting}{Verbatim}{commandchars=\\\{\}}
% Add ',fontsize=\small' for more characters per line
\usepackage{framed}
\definecolor{shadecolor}{RGB}{241,243,245}
\newenvironment{Shaded}{\begin{snugshade}}{\end{snugshade}}
\newcommand{\AlertTok}[1]{\textcolor[rgb]{0.68,0.00,0.00}{#1}}
\newcommand{\AnnotationTok}[1]{\textcolor[rgb]{0.37,0.37,0.37}{#1}}
\newcommand{\AttributeTok}[1]{\textcolor[rgb]{0.40,0.45,0.13}{#1}}
\newcommand{\BaseNTok}[1]{\textcolor[rgb]{0.68,0.00,0.00}{#1}}
\newcommand{\BuiltInTok}[1]{\textcolor[rgb]{0.00,0.23,0.31}{#1}}
\newcommand{\CharTok}[1]{\textcolor[rgb]{0.13,0.47,0.30}{#1}}
\newcommand{\CommentTok}[1]{\textcolor[rgb]{0.37,0.37,0.37}{#1}}
\newcommand{\CommentVarTok}[1]{\textcolor[rgb]{0.37,0.37,0.37}{\textit{#1}}}
\newcommand{\ConstantTok}[1]{\textcolor[rgb]{0.56,0.35,0.01}{#1}}
\newcommand{\ControlFlowTok}[1]{\textcolor[rgb]{0.00,0.23,0.31}{\textbf{#1}}}
\newcommand{\DataTypeTok}[1]{\textcolor[rgb]{0.68,0.00,0.00}{#1}}
\newcommand{\DecValTok}[1]{\textcolor[rgb]{0.68,0.00,0.00}{#1}}
\newcommand{\DocumentationTok}[1]{\textcolor[rgb]{0.37,0.37,0.37}{\textit{#1}}}
\newcommand{\ErrorTok}[1]{\textcolor[rgb]{0.68,0.00,0.00}{#1}}
\newcommand{\ExtensionTok}[1]{\textcolor[rgb]{0.00,0.23,0.31}{#1}}
\newcommand{\FloatTok}[1]{\textcolor[rgb]{0.68,0.00,0.00}{#1}}
\newcommand{\FunctionTok}[1]{\textcolor[rgb]{0.28,0.35,0.67}{#1}}
\newcommand{\ImportTok}[1]{\textcolor[rgb]{0.00,0.46,0.62}{#1}}
\newcommand{\InformationTok}[1]{\textcolor[rgb]{0.37,0.37,0.37}{#1}}
\newcommand{\KeywordTok}[1]{\textcolor[rgb]{0.00,0.23,0.31}{\textbf{#1}}}
\newcommand{\NormalTok}[1]{\textcolor[rgb]{0.00,0.23,0.31}{#1}}
\newcommand{\OperatorTok}[1]{\textcolor[rgb]{0.37,0.37,0.37}{#1}}
\newcommand{\OtherTok}[1]{\textcolor[rgb]{0.00,0.23,0.31}{#1}}
\newcommand{\PreprocessorTok}[1]{\textcolor[rgb]{0.68,0.00,0.00}{#1}}
\newcommand{\RegionMarkerTok}[1]{\textcolor[rgb]{0.00,0.23,0.31}{#1}}
\newcommand{\SpecialCharTok}[1]{\textcolor[rgb]{0.37,0.37,0.37}{#1}}
\newcommand{\SpecialStringTok}[1]{\textcolor[rgb]{0.13,0.47,0.30}{#1}}
\newcommand{\StringTok}[1]{\textcolor[rgb]{0.13,0.47,0.30}{#1}}
\newcommand{\VariableTok}[1]{\textcolor[rgb]{0.07,0.07,0.07}{#1}}
\newcommand{\VerbatimStringTok}[1]{\textcolor[rgb]{0.13,0.47,0.30}{#1}}
\newcommand{\WarningTok}[1]{\textcolor[rgb]{0.37,0.37,0.37}{\textit{#1}}}

\usepackage{longtable,booktabs,array}
\usepackage{calc} % for calculating minipage widths
% Correct order of tables after \paragraph or \subparagraph
\usepackage{etoolbox}
\makeatletter
\patchcmd\longtable{\par}{\if@noskipsec\mbox{}\fi\par}{}{}
\makeatother
% Allow footnotes in longtable head/foot
\IfFileExists{footnotehyper.sty}{\usepackage{footnotehyper}}{\usepackage{footnote}}
\makesavenoteenv{longtable}
\usepackage{graphicx}
\makeatletter
\newsavebox\pandoc@box
\newcommand*\pandocbounded[1]{% scales image to fit in text height/width
  \sbox\pandoc@box{#1}%
  \Gscale@div\@tempa{\textheight}{\dimexpr\ht\pandoc@box+\dp\pandoc@box\relax}%
  \Gscale@div\@tempb{\linewidth}{\wd\pandoc@box}%
  \ifdim\@tempb\p@<\@tempa\p@\let\@tempa\@tempb\fi% select the smaller of both
  \ifdim\@tempa\p@<\p@\scalebox{\@tempa}{\usebox\pandoc@box}%
  \else\usebox{\pandoc@box}%
  \fi%
}
% Set default figure placement to htbp
\def\fps@figure{htbp}
\makeatother





\setlength{\emergencystretch}{3em} % prevent overfull lines

\providecommand{\tightlist}{%
  \setlength{\itemsep}{0pt}\setlength{\parskip}{0pt}}



 


% Paquetes básicos
\usepackage{booktabs}
\usepackage{pdfpages}
\usepackage{amsthm}
\usepackage{graphicx}

% Configuración de márgenes y espaciado
\usepackage{geometry}
\geometry{
  a4paper,
  top=2.5cm,
  bottom=2.5cm,
  left=3cm,
  right=3cm
}

% Espaciado de teoremas
\makeatletter
\def\thm@space@setup{%
  \thm@preskip=8pt plus 2pt minus 4pt
  \thm@postskip=\thm@preskip
}
\makeatother

% Configuración de portada
\let\oldmaketitle\maketitle
\AtBeginDocument{\let\maketitle\relax}

% Mejoras de hipervínculos
\usepackage{hyperref}
\hypersetup{
  colorlinks=true,
  linkcolor=blue,
  filecolor=magenta,
  urlcolor=cyan,
  pdfpagemode=FullScreen,
}

% Configuración de listados de código
\usepackage{listings}
\lstset{
  basicstyle=\ttfamily\small,
  breaklines=true,
  frame=single,
  numbers=left,
  numberstyle=\tiny,
  captionpos=b
}

% NOTA: Los entornos filus y aisa están definidos en fonts.tex
% NO duplicar definiciones aquí
% ============================================
% CONFIGURACIÓN DE FUENTES CON FONTSPEC
% ============================================
\usepackage{fontspec}
\usepackage{xcolor}

% Fuente principal (texto general)
\setmainfont{LibertinusSerif}[
  Path=_fonts/,
  Extension=.otf,
  UprightFont=*-Regular,
  BoldFont=*-Bold,
  ItalicFont=*-Italic,
  BoldItalicFont=*-BoldItalic
]

% Fuente sans-serif (títulos)
\setsansfont{LibertinusSans}[
  Path=_fonts/,
  Extension=.otf,
  UprightFont=*-Regular,
  BoldFont=*-Bold,
  ItalicFont=*-Italic
]

% Fuente monoespaciada (código)
\setmonofont{FiraCode}[
  Path=_fonts/,
  Extension=.ttf,
  UprightFont=*-Regular,
  Scale=0.85
]

% ============================================
% TIPOGRAFÍAS PARA LOS PERSONAJES
% ============================================

% Tipografía para Filus (el humano) - Serif elegante
\newfontfamily\filusfont{LibertinusSerif}[
  Path=_fonts/,
  Extension=.otf,
  UprightFont=*-Regular,
  BoldFont=*-Bold,
  ItalicFont=*-Italic,
  Color=2C3E50
]

% Tipografía para Aisa (la IA) - Sans-serif moderna
\newfontfamily\aisafont{LibertinusSans}[
  Path=_fonts/,
  Extension=.otf,
  UprightFont=*-Regular,
  BoldFont=*-Bold,
  ItalicFont=*-Italic,
  Color=2980B9
]

% ============================================
% COMANDOS PERSONALIZADOS PARA DIÁLOGOS
% ============================================

% Comandos cortos (usando nombres únicos)
\newcommand{\filustext}[1]{%
  \begin{quote}
    {\filusfont\textbf{Filus:} #1}
  \end{quote}
}

\newcommand{\aisatext}[1]{%
  \begin{quote}
    {\aisafont\textbf{Aisa:} #1}
  \end{quote}
}

% Entornos largos (usando nombres únicos)
\newenvironment{dialogofilus}{%
  \begin{quote}
  \filusfont
  \noindent\textbf{Filus:}\par
}{%
  \end{quote}
}

\newenvironment{dialogoaisa}{%
  \begin{quote}
  \aisafont
  \noindent\textbf{Aisa:}\par
}{%
  \end{quote}
}

% ============================================
% PERSONALIZACIÓN DE TÍTULOS CON KOMA-SCRIPT
% ============================================

\addtokomafont{chapter}{\sffamily\bfseries}
\addtokomafont{section}{\sffamily\bfseries}
\addtokomafont{subsection}{\sffamily\bfseries}

% ============================================
% COLORES
% ============================================

\definecolor{filuscolor}{HTML}{2C3E50}
\definecolor{aisacolor}{HTML}{2980B9}
\definecolor{codecolor}{HTML}{3E4451}
\makeatletter
\@ifpackageloaded{tcolorbox}{}{\usepackage[skins,breakable]{tcolorbox}}
\@ifpackageloaded{fontawesome5}{}{\usepackage{fontawesome5}}
\definecolor{quarto-callout-color}{HTML}{909090}
\definecolor{quarto-callout-note-color}{HTML}{0758E5}
\definecolor{quarto-callout-important-color}{HTML}{CC1914}
\definecolor{quarto-callout-warning-color}{HTML}{EB9113}
\definecolor{quarto-callout-tip-color}{HTML}{00A047}
\definecolor{quarto-callout-caution-color}{HTML}{FC5300}
\definecolor{quarto-callout-color-frame}{HTML}{acacac}
\definecolor{quarto-callout-note-color-frame}{HTML}{4582ec}
\definecolor{quarto-callout-important-color-frame}{HTML}{d9534f}
\definecolor{quarto-callout-warning-color-frame}{HTML}{f0ad4e}
\definecolor{quarto-callout-tip-color-frame}{HTML}{02b875}
\definecolor{quarto-callout-caution-color-frame}{HTML}{fd7e14}
\makeatother
\makeatletter
\@ifpackageloaded{bookmark}{}{\usepackage{bookmark}}
\makeatother
\makeatletter
\@ifpackageloaded{caption}{}{\usepackage{caption}}
\AtBeginDocument{%
\ifdefined\contentsname
  \renewcommand*\contentsname{Table of contents}
\else
  \newcommand\contentsname{Table of contents}
\fi
\ifdefined\listfigurename
  \renewcommand*\listfigurename{List of Figures}
\else
  \newcommand\listfigurename{List of Figures}
\fi
\ifdefined\listtablename
  \renewcommand*\listtablename{List of Tables}
\else
  \newcommand\listtablename{List of Tables}
\fi
\ifdefined\figurename
  \renewcommand*\figurename{Figure}
\else
  \newcommand\figurename{Figure}
\fi
\ifdefined\tablename
  \renewcommand*\tablename{Table}
\else
  \newcommand\tablename{Table}
\fi
}
\@ifpackageloaded{float}{}{\usepackage{float}}
\floatstyle{ruled}
\@ifundefined{c@chapter}{\newfloat{codelisting}{h}{lop}}{\newfloat{codelisting}{h}{lop}[chapter]}
\floatname{codelisting}{Listing}
\newcommand*\listoflistings{\listof{codelisting}{List of Listings}}
\makeatother
\makeatletter
\makeatother
\makeatletter
\@ifpackageloaded{caption}{}{\usepackage{caption}}
\@ifpackageloaded{subcaption}{}{\usepackage{subcaption}}
\makeatother
\usepackage{bookmark}
\IfFileExists{xurl.sty}{\usepackage{xurl}}{} % add URL line breaks if available
\urlstyle{same}
\hypersetup{
  pdftitle={El Arquitecto Invisible},
  pdfauthor={Filus \& Aisa},
  hidelinks,
  pdfcreator={LaTeX via pandoc}}


\title{El Arquitecto Invisible}
\author{Filus \& Aisa}
\date{2025-10-05}
\begin{document}
\frontmatter
\maketitle

\includepdf[pages={1}, scale=1]{img/cover_page_scriberia.pdf}
\newpage

\let\maketitle\oldmaketitle
\maketitle

\renewcommand*\contentsname{Table of contents}
{
\setcounter{tocdepth}{2}
\tableofcontents
}

\mainmatter
\bookmarksetup{startatroot}

\chapter*{Acknowledgement}\label{acknowledgement}
\addcontentsline{toc}{chapter}{Acknowledgement}

\markboth{Acknowledgement}{Acknowledgement}

\begin{figure}[H]

{\centering \includegraphics[width=5.72917in,height=4.16667in]{img/r-dev-guide-without-text.jpg}

}

\caption{R Development Guide. This illustration is created by Scriberia
with The Turing Way community, used under a CC-BY 4.0 licence. DOI:
https://zenodo.org/records/13882307}

\end{figure}%

This guide draws on documentation and articles written by the R Core
Team. The first version of the guide was heavily influenced by the
\href{https://devguide.python.org/}{Python Developer's Guide}.

Initial chapters of the guide were developed by Saranjeet Kaur Bhogal,
in a project funded by the R Foundation, mentored by Heather Turner and
Michael Lawrence. This initial version was upgraded in a
\href{https://github.com/rstats-gsod/gsod2022/wiki/GSOD-2022-Proposal}{Google
Season of Docs 2022} project with Saranjeet Kaur Bhogal and Lluís
Revilla Sancho working as technical writers managed by Nicolas Bennett
and overseen by a
\href{https://github.com/rstats-gsod/gsod2022/wiki/GSOD-2022-Proposal\#steering-committee}{Steering
Committee} including representatives from R Core and the wider R
community.

This guide has benefited and continues to benefit from varied
contributions by several
\href{https://github.com/r-devel/rdevguide\#contributors-}{contributors}.

\begin{figure}[H]

{\centering \pandocbounded{\includegraphics[keepaspectratio]{img/ccby.png}}

}

\caption{License: CC BY 4.0}

\end{figure}%

This project is licensed under a
\href{https://creativecommons.org/licenses/by/4.0/}{Creative Commons
Attribution 4.0 International (CC BY 4.0)}. Some pages may contain
materials that are subject to copyright, in which case you will see the
copyright notice.

\bookmarksetup{startatroot}

\chapter{Bienvenidos a El Arquitecto
Invisible}\label{bienvenidos-a-el-arquitecto-invisible}

Este es el texto principal del capítulo, que usará la fuente por defecto
del tema \texttt{cosmo}.

\section{Una sección con estilos
personalizados}\label{una-secciuxf3n-con-estilos-personalizados}

Podemos aplicar estilos específicos a diferentes partes del texto.

\begin{dialogofilus}
Este párrafo está escrito con la fuente Cambria, aplicando la clase
\texttt{.filus}. Es ideal para citas o fragmentos que necesiten un toque
más clásico y editorial.
\end{dialogofilus}

\begin{dialogoaisa}
Este otro párrafo utiliza la fuente Segoe UI, gracias a la clase
\texttt{.aisa}. Esta fuente es moderna y muy legible, perfecta para
notas o comentarios adicionales.
\end{dialogoaisa}

Y aquí tenemos un ejemplo de código:

\begin{Shaded}
\begin{Highlighting}[]
\KeywordTok{function} \FunctionTok{helloWorld}\NormalTok{() \{}
  \CommentTok{// Este código usará la fuente Source Code Pro}
  \BuiltInTok{console}\OperatorTok{.}\FunctionTok{log}\NormalTok{(}\StringTok{"¡Hola, mundo tipográfico!"}\NormalTok{)}\OperatorTok{;}
\NormalTok{\}}
\end{Highlighting}
\end{Shaded}

\bookmarksetup{startatroot}

\chapter{Capítulo 1}\label{capuxedtulo-1}

Este es texto normal.

\filustext{Hola Aisa, ¿cómo funcionan las tipografías?}

\aisatext{Las tipografías nos dan personalidad visual. Yo uso sans-serif porque soy una IA moderna, tú usas serif porque eres humano clásico.}

\section{Código de ejemplo}\label{cuxf3digo-de-ejemplo}

\begin{Shaded}
\begin{Highlighting}[]
\FunctionTok{print}\NormalTok{(}\StringTok{"Hola mundo filus lua"}\NormalTok{)}
\end{Highlighting}
\end{Shaded}

\section{Diálogos más extensos}\label{diuxe1logos-muxe1s-extensos}

Para conversaciones más largas, ahora puedes usar la sintaxis unificada:

\begin{dialogofilus}
¿Y qué pasa si quiero desarrollar una idea más compleja? ¿Puedo usar
múltiples párrafos?

Porque a veces necesito expresar pensamientos que requieren más espacio
y estructura.
\end{dialogofilus}

\begin{dialogoaisa}
Por supuesto, Filus. Los bloques Div están diseñados precisamente para
eso.

Puedes escribir tantos párrafos como necesites, y todo mantendrá tu
tipografía distintiva. Y ahora funciona tanto en HTML como en PDF
gracias al filtro Lua.
\end{dialogoaisa}

\section{Continuando con texto
normal}\label{continuando-con-texto-normal}

Después de los diálogos, el texto narrativo vuelve a la tipografía
principal del libro.

\bookmarksetup{startatroot}

\chapter{Capítulo tu}\label{capuxedtulo-tu}

Este es texto normal.

\begin{Shaded}
\begin{Highlighting}[]
\FunctionTok{print}\NormalTok{(}\StringTok{"Hola mundo"}\NormalTok{)}
\end{Highlighting}
\end{Shaded}

\section{Diálogos más extensos}\label{diuxe1logos-muxe1s-extensos-1}

Para conversaciones más largas, puedes usar los entornos:

\begin{dialogofilus}
¿Y qué pasa si quiero desarrollar una idea más compleja? ¿Puedo usar
múltiples párrafos?

Porque a veces necesito expresar pensamientos que requieren más espacio
y estructura.
\end{dialogofilus}

\begin{dialogoaisa}
¿Y qué pasa si quiero desarrollar una idea más compleja? ¿Puedo usar
múltiples párrafos?

Porque a veces necesito expresar pensamientos que requieren más espacio
y estructura. Por supuesto, Filus. Los entornos largos están diseñados
precisamente para eso.

Puedes escribir tantos párrafos como necesites, y todo mantendrá tu
tipografía distintiva.
\end{dialogoaisa}

\begin{verbatim}

## Continuando con texto normal

Después de los diálogos, el texto narrativo vuelve
a la tipografía principal del libro.

`<!-- quarto-file-metadata: eyJyZXNvdXJjZURpciI6Ii4ifQ== -->`{=html}

```{=html}
<!-- quarto-file-metadata: eyJyZXNvdXJjZURpciI6Ii4iLCJib29rSXRlbVR5cGUiOiJjaGFwdGVyIiwiYm9va0l0ZW1OdW1iZXIiOjQsImJvb2tJdGVtRmlsZSI6ImZpbHVzLnFtZCIsImJvb2tJdGVtRGVwdGgiOjB9 -->
\end{verbatim}

\bookmarksetup{startatroot}

\chapter{}\label{section}

\begin{tcolorbox}[enhanced jigsaw, leftrule=.75mm, toprule=.15mm, colback=white, arc=.35mm, colframe=quarto-callout-color-frame, breakable, rightrule=.15mm, bottomrule=.15mm, left=2mm, opacityback=0]

Hi i am Filus

\end{tcolorbox}

\begin{center}\rule{0.5\linewidth}{0.5pt}\end{center}

\[ v := vec(x_1, x_2, x_3) \]

\begin{center}\rule{0.5\linewidth}{0.5pt}\end{center}

\[ (x+y)^n = \sum_{k=0}^n {n \choose k}x^{n-k}y^k = \sum_{k=0}^n {n \choose k}x^{k}y^{n-k} \]

\bookmarksetup{startatroot}

\chapter{Summary}\label{summary}

In summary, this book has no content whatsoever.

\bookmarksetup{startatroot}

\chapter*{References}\label{references}
\addcontentsline{toc}{chapter}{References}

\markboth{References}{References}

\phantomsection\label{refs}


\backmatter


\end{document}
